\section{Conclusion}
\label{sec:conclusion}

This paper has set out a discrete model of reality built entirely from experience. From one axiom, that to exist is to experience or be experienced, and two constraints, that the base is finite and that the only entity is the vibe, the model derives a working account of its own dynamics, the energy experience minimizes, the discrete time it runs in, the necessity and character of its own beginning, the unbounded growth that follows, the emergence of selves as self-maintaining patterns, the emergence of space and its hyperbolic shape, the mechanics of how an agent moves through a substrate woven from other agents, and a route toward building matter without any new primitive.

Two commitments give the model its shape. Discreteness keeps it finite and simulable, so every claim becomes a computational experiment and every quantity an exact integer rather than an approximation. Monism keeps it built from one kind of thing, so links, fields, particles, space, and time are all patterns of vibes noting vibes, and the inside of a particle is a small system of vibes rather than a new substance. Read together, these two demands turn a philosophical intuition into a constructive program.

The single thesis is that experience is the foundation, not an emergent extra. The discrete mathematics that describes how the substrate updates is a description of its form, layered on top, in the way that the equations of a fluid describe molecular motion without being what the fluid is made of. Physics is recovered from experience. Experience is not recovered from physics.

The model remains a philosophical and mathematical proposal rather than a tested theory. Its connection to established physics is a program, its account of matter is a frontier, and its deepest cosmological questions are openly unresolved. What it offers is a coherent, discrete, experience-first ontology with a concrete build target, stated precisely enough to be simulated, criticized, and refined.
